\section{Evaluation}\label{sec:evaluation}

\begin{table}[!t]
  \centering
  \caption{Offline results: \mossinstruct{} (0708 release) against
  open-source models of comparable scale, blocked by capability domain.
  Bold = best, underline = second best, ``--'' = not reported;
  LLaVA-OV-2 = LLaVA-OneVision-2-8B. Evaluation protocol and baseline
  provenance are spelled out in \S\ref{sec:eval-offline}.}
  \label{tab:offline}
  \footnotesize
\setlength{\tabcolsep}{4.5pt}
\renewcommand{\arraystretch}{1.12}
\begin{tabular}{@{}lccccc@{}}
\toprule
Benchmark & \makecell{\textbf{MOSS-VL}\\\textbf{(ours)}} & \makecell{Qwen3-VL\\8B} & \makecell{Qwen2.5-VL\\7B} & \makecell{LLaVA-OV-2\\8B} & \makecell{Gemma-4\\12B-IT} \\
\midrule
\multicolumn{6}{@{}l}{\itshape Multimodal perception} \\
MMBench-EN (v1.1)~\citep{Liu2023mmbench} & \bnum{88.1} & 84.8 & 83.2 & \snum{85.8} & 82.7 \\
MMStar~\citep{Chen2024right} & 66.0 & \snum{70.9} & 63.9 & 64.8 & \bnum{74.9} \\
RealWorldQA       & 68.0 & \bnum{71.5} & 68.5 & \snum{69.7} & 65.6 \\
MME-RealWorld~\citep{Zhang2024realworld} & \bnum{66.3} & -- & \snum{57.4} & -- & 46.9 \\
BLINK~\citep{Fu2024blink} & \bnum{78.0} & \snum{69.1} & 56.4 & 63.5 & 65.6 \\
POPE~\citep{Li2023evaluating} & \bnum{89.4} & -- & \snum{87.4} & -- & 81.4 \\
MMMU (val)~\citep{Yue2023mmmu} & 51.1 & \snum{69.6} & 58.6 & -- & \bnum{69.7} \\
CountBench~\citep{Paiss2023teaching} & 85.9 & 80.5 & -- & \snum{89.0} & \bnum{90.2} \\
CVBench~\citep{Tong2024cambrian} & 85.7 & \snum{86.7} & -- & \bnum{87.7} & 84.6 \\
V*~\citep{Wu2023guided} & \bnum{89.0} & 85.3 & -- & \snum{85.9} & 51.8 \\
MuirBench~\citep{Wang2024muirbench} & 39.9 & \bnum{64.4} & 59.6 & -- & \snum{61.6} \\
AI2D~\citep{Kembhavi2016diagram} & 81.9 & \bnum{85.7} & 83.9 & 84.3 & \snum{84.8} \\
\midrule
\multicolumn{6}{@{}l}{\itshape Video understanding} \\
VideoMME~\citep{Fu2024video} & 68.1 & \snum{71.4} & 65.1 & \bnum{71.9} & 60.5 \\
VideoMME-v2~\citep{Fu2026video} & \bnum{12.7} & \snum{12.4} & 10.3 & -- & -- \\
VideoMME-v2 (sub)~\citep{Fu2026video} & \snum{16.5} & \bnum{18.2} & -- & -- & -- \\
VideoMME-Logical~\citep{Kwan2026video} & \bnum{17.1} & \snum{11.9} & 7.4 & -- & 10.8 \\
MLVU (dev)~\citep{Zhou2024mlvu} & \snum{76.8} & \bnum{78.1} & 70.2 & 76.6 & -- \\
LongVideoBench~\citep{Wu2024longvideobench} & 65.9 & \bnum{68.0} & 56.0 & \snum{66.9} & 58.2 \\
LVBench~\citep{Wang2024lvbench} & 51.1 & \bnum{58.0} & 45.3 & \snum{55.5} & 37.3 \\
EgoSchema (sub)~\citep{Mangalam2023egoschema} & \bnum{67.0} & -- & \snum{65.0} & -- & 62.2 \\
MVBench~\citep{Li2023mvbench} & 66.7 & \snum{68.7} & \bnum{69.6} & 66.2 & -- \\
VSI-Bench~\citep{Yang2024thinking} & \snum{62.2} & 59.4 & 28.3 & \bnum{70.9} & 25.9 \\
Minerva~\citep{Nagrani2025minerva} & \bnum{40.5} & -- & -- & -- & \snum{32.3} \\
TimeLens-Charades~\citep{Zhang2025timelens} & 51.5 & \bnum{56.0} & 43.6 & \snum{53.5} & -- \\
TimeLens-ANet~\citep{Zhang2025timelens} & \snum{49.1} & 46.8 & 31.4 & \bnum{53.8} & -- \\
TimeLens-QVH~\citep{Zhang2025timelens} & \snum{60.0} & 59.4 & 31.6 & \bnum{66.4} & -- \\
TOMATO~\citep{Shangguan2024tomato} & \bnum{39.5} & \snum{34.6} & -- & -- & 31.9 \\
\midrule
\multicolumn{6}{@{}l}{\itshape Grounding} \\
RefCOCO-REC~\citep{Yu2016modeling} & 84.4 & \bnum{91.6} & \snum{90.0} & -- & -- \\
Ref-Adv~\citep{Akula2020words} & \bnum{57.0} & 47.2 & \snum{49.3} & -- & -- \\
\midrule
\multicolumn{6}{@{}l}{\itshape Document / OCR} \\
DocVQA (val)~\citep{Mathew2020docvqa} & 89.6 & \bnum{96.1} & \snum{95.7} & 95.2 & 80.8 \\
ChartQA~\citep{Masry2022chartqa} & \snum{87.8} & \bnum{89.6} & 87.3 & 85.9 & 51.2 \\
InfoVQA (val)~\citep{Mathew2021infographicvqa} & 68.9 & \bnum{83.4} & \snum{82.6} & 74.4 & 52.0 \\
OCRBench~\citep{Liu2023ocrbench} & 86.1 & \bnum{89.6} & \snum{86.4} & 78.2 & 76.9 \\
OCRBench-v2~\citep{Fu2024ocrbench} & \snum{57.4} & \bnum{63.3} & 56.8 & -- & 39.6 \\
OmniDocBench (v1.6)~\citep{Ouyang2024omnidocbench} & \bnum{88.9} & \snum{84.9} & -- & -- & -- \\
\midrule
\multicolumn{6}{@{}l}{\itshape Reasoning} \\
VLMsAreBlind~\citep{Rahmanzadehgervi2024vision} & 63.5 & \snum{74.0} & -- & -- & \bnum{75.9} \\
VisuLogic~\citep{Xu2025visulogic} & \bnum{27.5} & 22.5 & \snum{26.0} & -- & -- \\
ERQA~\citep{Team2025gemini} & \bnum{45.8} & \bnum{45.8} & -- & \snum{43.3} & 40.8 \\
EmbSpatial~\citep{Du2024embspatial} & 70.7 & \bnum{78.5} & -- & \snum{78.1} & 72.5 \\
\bottomrule
\end{tabular}

\end{table}

\begin{figure}[t]
  \centering
  \includegraphics[width=\linewidth]{fig-efficiency.pdf}
  \caption{Measured serving latency of \mossvl{} vs.\ Qwen3-VL-8B --- the
  same Qwen3-8B language backbone, isolating the vision-integration
  architecture. Both models serve offline with SGLang on a single H200
  (TP=1, BF16, identical engine version) with an identical
  generation-length cap, which every run reaches; every point is the mean
  of five independent cold starts (error bars: sample std). (a,\,c) With ViT
  output matched, the time-to-first-token (TTFT) gap widens from $2.8\times$ to
  $5.1\times$ as visual context grows, and end-to-end latency from
  $1.9\times$ to $4.3\times$. (b,\,d) On the same video
  at the same resolution and frame count: Qwen3-VL compresses the
  temporal axis $2\times$, whereas \mossvl{} forgoes temporal compression
  so that each arriving frame can be encoded immediately --- carrying
  about twice the vision tokens on the same input --- yet it never
  falls behind, and its lead grows with stream length. Slower latency
  growth follows from the append-only
  cross-attention design: visual tokens never enter the decoded sequence
  (\S\ref{sec:architecture}).}
  \label{fig:efficiency}
\end{figure}

We evaluate each model in the regime it is built for: \mossinstruct{}
on an offline suite of 39 benchmarks across five capability domains
(Table~\ref{tab:offline}), and \mossrealtime{} on four streaming
benchmarks---OVO-Bench \citep{Li2025bench}, OmniMMI
\citep{Wang2025omnimmi}, StreamingBench \citep{Lin2024streamingbench},
and ProactiveVideoQA \citep{Wang2025proactivevideoqa}---which together
cover levels L2--L4 of the capability hierarchy in
Table~\ref{tab:levels}. All streaming evaluation is carried out in
streaming fashion: frames are fed as they would arrive, and every
model runs under its own streaming protocol. Measured serving
efficiency (\S\ref{sec:eval-efficiency}) and qualitative real-time
sessions (\S\ref{sec:eval-showcase}) complete the picture.

\subsection{Offline Results}\label{sec:eval-offline}

Table~\ref{tab:offline} compares \mossinstruct{} with open models of
comparable scale---Qwen3-VL-8B \citep{Bai2025qwen}, Qwen2.5-VL-7B
\citep{Bai2025qwenvl25}, LLaVA-OneVision-2-8B \citep{An2026llava}, and
Gemma-4-12B-IT \citep{Team2026gemma}---over multimodal perception, video
understanding, grounding, document/OCR, and reasoning.
\mossvl{} results sample video at 1\,fps with at most 768 frames,
following a benchmark's official protocol wherever one is prescribed;
DocVQA and InfoVQA use the validation split. Baseline numbers are
taken from the respective official reports and reflect their authors'
inference settings, which may differ from ours, particularly in video
frame count. The exceptions are the Gemma-4-12B-IT column, which we
evaluated ourselves under the same protocol as \mossvl{}, and the
Qwen3-VL entry on OmniDocBench (v1.6), evaluated with the Markdown
prompt from the official Qwen3-VL cookbook; the remaining OmniDocBench
baselines are omitted, as their official numbers are not
metric-comparable.

Perception is the strongest block: \mossinstruct{} takes five of the
twelve rows---MMBench-EN (88.1), POPE (89.4), V* (89.0), and both
MME-RealWorld (66.3) and BLINK (78.0) by 8.9-point margins. On video,
it leads the temporal-reasoning sets---Minerva (40.5),
TOMATO (39.5), and VideoMME-Logical (17.1), each by 4.9 points or
more---plus EgoSchema (67.0), consistent with the
perception and temporal-understanding foundations built in
\S\ref{sec:pretraining}. The wins extend across the remaining
domains: the adversarial Ref-Adv grounding set goes to \mossinstruct{}
by 7.7 points (57.0), OmniDocBench document parsing by 4.0 (88.9), and on
reasoning it takes VisuLogic (27.5) and shares the top ERQA score
(45.8). The main gaps sit in MMMU, document understanding, and
standard grounding, both referring (RefCOCO-REC) and temporal
(TimeLens); the first two we return to in \S\ref{sec:discussion}.

\subsection{Streaming Benchmarks}\label{sec:eval-streaming}

Table~\ref{tab:streaming} reports subset-level results on the four
streaming benchmarks. Baselines are open-source streaming
models---AURA \citep{Lu2026aura}, M4 (released with OmniMMI itself)
\citep{Wang2025omnimmi}, ROMA \citep{Tian2026roma},
JoyAI-VL-Interaction \citep{Yao2026joyai}, VideoChat3-4B
\citep{Li2026videochat3}, ViSpeak-7B \citep{Fu2025vispeak}, and the
MMDuet family
\citep{Wang2024videollm,Wang2025proactivevideoqa,Wang2025mmduet2}; the
panel follows each benchmark's published coverage and therefore differs
across the four. Baseline numbers are taken from the respective
official reports, except MMDuet's OmniMMI entry, which its own report
does not cover and is taken from the ROMA report; \mossrealtime{}
numbers are from our own evaluation under each benchmark's official
protocol. %
\mossrealtime{} runs an 8.2B language backbone, comparable to
the 7--8B-class baselines; its additional parameters lie outside the
decoded sequence, in the vision encoder and the cross-attention stack
(\S\ref{sec:architecture}).

\begin{table}[!t]
  \centering
  \caption{Streaming benchmark results: \mossrealtime{} against
  open-source streaming baselines at subset level. Bold = best,
  underline = second best; ``--'' = subset not reported in that
  model's official report. Subset abbreviations, each benchmark's Avg
  convention, and data provenance are spelled out in
  \S\ref{sec:eval-streaming}.}
  \label{tab:streaming}
  \footnotesize
\setlength{\tabcolsep}{6pt}
\renewcommand{\arraystretch}{1.12}
\begin{tabular}{@{}lccccc@{}}
\toprule
\multicolumn{6}{@{}l}{\itshape OVO-Bench~\citep{Li2025bench}} \\
Subset & \textbf{MOSS-VL-Realtime} & AURA & JoyAI-VL-Interaction & VideoChat3-4B & ViSpeak-7B \\
\midrule
FAR  & \bnum{62.1} & \snum{55.8} & -- & -- & 54.2 \\
BT   & \bnum{72.6} & \snum{60.4} & -- & -- & 57.5 \\
RTVP & \snum{75.9} & \bnum{79.8} & -- & -- & 66.3 \\
Avg  & \bnum{70.2} & \snum{65.3} & 59.2 & 57.8 & 59.3 \\
\midrule
\multicolumn{6}{@{}l}{\itshape OmniMMI~\citep{Wang2025omnimmi}} \\
Subset & \textbf{MOSS-VL-Realtime} & AURA & M4 & ROMA & MMDuet \\
\midrule
PA  & \bnum{66.0} & \snum{37.5} & 25.5 & \snum{37.5} & 22.0 \\
SG  & \snum{21.7} & \bnum{24.0} & 5.7 & -- & -- \\
MD  & \bnum{10.7} & \snum{7.7} & 1.7 & -- & -- \\
AP  & \bnum{33.5} & \snum{32.0} & \bnum{33.5} & -- & -- \\
SI  & \bnum{31.5} & \snum{26.0} & 9.0 & -- & -- \\
Avg & \bnum{32.7} & \snum{25.4} & 15.1 & -- & -- \\
\midrule
\multicolumn{6}{@{}l}{\itshape StreamingBench~\citep{Lin2024streamingbench}} \\
Subset & \textbf{MOSS-VL-Realtime} & AURA & VideoChat3-4B & ViSpeak-7B & \\
\midrule
PO  & \bnum{60.0} & \snum{53.2} & -- & 50.8 & \\
RT  & 82.9 & \bnum{83.2} & \snum{83.0} & 70.4 & \\
CTX & \snum{56.4} & \bnum{59.0} & -- & 43.9 & \\
SQA & \snum{50.4} & \bnum{57.2} & -- & 39.2 & \\
Avg (visual) & \snum{69.7} & \bnum{71.1} & -- & 57.2 & \\
\midrule
\multicolumn{6}{@{}l}{\itshape ProactiveVideoQA~\citep{Wang2025proactivevideoqa}} \\
Subset & \textbf{MOSS-VL-Realtime} & MMDuet2-RL & MMDuet+rm & MMDuet & VideoChat3-4B \\
\midrule
WEB & \bnum{55.4} & \snum{53.3} & 43.5 & 38.9 & 38.4 \\
EGO & \snum{47.8} & 33.6 & \bnum{52.2} & 46.0 & 28.1 \\
TV  & \bnum{50.2} & \snum{43.4} & 32.6 & 21.1 & 34.7 \\
VAD & \snum{35.3} & 28.9 & \bnum{42.5} & 27.4 & 25.1 \\
Avg & \bnum{47.2} & 39.8 & \snum{42.7} & 33.4 & 31.6 \\
\bottomrule
\end{tabular}

\end{table}

Each benchmark reports at subset level. OVO-Bench separates forward
active responding (FAR), backward tracing (BT), and real-time visual
perception (RTVP). OmniMMI covers proactive alerting (PA), dynamic
state grounding (SG), multi-turn dependency (MD), action prediction
(AP), and speaker identification (SI); its Avg is the mean of the five
subsets, shown only for models with all five reported. StreamingBench
groups real-time visual understanding (RT) and contextual
understanding (CTX), with proactive output (PO) and sequential QA
(SQA), subsets of CTX, listed separately for the proactive analysis;
Avg (visual) is the mean of the RT and CTX group scores---\mossvl{}
takes no audio input, so the audio-dependent Omni-Source group is not
evaluated and no official overall score is reported. ProactiveVideoQA
splits by video source into web (WEB), egocentric (EGO), and TV-series
(TV) video QA, plus video anomaly detection (VAD).

\mossrealtime{} posts the best average on three of the four
benchmarks---OVO-Bench (70.2 vs.\ 65.3 for the runner-up), OmniMMI
(32.7 vs.\ 25.4), and ProactiveVideoQA (47.2 vs.\ 42.7)---and is
second on StreamingBench's visual average (69.7 vs.\ AURA's 71.1).

The subset pattern says more than the averages. The three subsets that
squarely test proactive behavior---speaking unprompted, at the right
moment---all go to \mossrealtime{}: Proactive Alerting on OmniMMI
(66.0 vs.\ 37.5), Proactive Output on StreamingBench (60.0 vs.\ 53.2),
and Forward Active Responding on OVO-Bench (62.1 vs.\ 55.8).
ProactiveVideoQA, proactive in every subset, follows in aggregate, and
Backward Tracing on OVO-Bench (72.6 vs.\ 60.4) shows the accumulated
stream history staying usable. The wins concentrate where response
timing is the skill under test---the behavior Realtime-SFT
supervises directly (\S\ref{sec:realtime-sft}); where a subset reduces
to perception QA over the current scene, AURA keeps the edge.

\subsection{Inference Efficiency}\label{sec:eval-efficiency}

Figure~\ref{fig:efficiency} measures serving latency against
Qwen3-VL-8B, which shares the Qwen3-8B language backbone
\citep{Yang2025qwen}; the comparison therefore isolates the
vision-integration architecture. Both models run offline SGLang
serving on a single H200 (\S\ref{sec:infrastructure}). With ViT output
matched, the time-to-first-token gap widens from $2.8\times$ to
$5.1\times$ as visual context grows, and end-to-end latency from
$1.9\times$ to $4.3\times$. The same-video comparison is stricter for
us: \mossvl{} forgoes temporal compression so that each arriving frame
can be encoded immediately (\S\ref{sec:realtime-arch}), and thus
carries about twice the vision tokens of Qwen3-VL on identical
input---yet it never falls behind at any measured point. The advantage widens
with visual context, which is exactly the regime a real-time assistant
occupies: visual history accumulates by the minute while replies must
keep arriving on time.

\subsection{Qualitative Results}\label{sec:eval-showcase}

\begin{figure}[t]
  \centering
  \includegraphics[width=\linewidth]{fig-showcase.pdf}
  \caption{\mossrealtime{} in the wild: two screen-captured sessions from
  the released live demo, running on a single H200, shown as frame strips
  with excerpts of the model's outputs on a shared timeline (long outputs
  truncated with ``\ldots''; $t=0$ at the user instruction). (a) Under
  a standing conditional instruction (``say \emph{great!}\ whenever the
  cat touches the carrot, otherwise stay silent''), the model stays silent
  throughout and fires exactly at each of the four contacts; the hollow
  marker samples one of the silent spans, where the per-frame output is
  the silence token. (b) Under a
  single instruction (``commentate the video live''), commentary begins
  within a second and tracks a free-kick sequence --- set-up, the
  referee's whistle, the strike, the celebration, and the updated
  scoreline --- in a professional broadcast register. Outputs are in
  Chinese, the demo language; English translations of both sessions are
  given in Appendix~\ref{app:transcripts}, and these and further live
  cases can be viewed on the official blog
  (\url{https://openmoss.ai/MOSS-VL/}).}
  \label{fig:showcase}
\end{figure}

Figure~\ref{fig:showcase} shows the installed behaviors in live
operation, with two sessions from the released
demo. Under a standing conditional instruction, the model holds
silence over the full stream and fires at each of the four target
contacts, and only there; under a single commentary instruction, it
opens within a second and tracks a free-kick sequence through the
whistle, the strike, the celebration, and the updated scoreline.
Between them, the sessions exercise the core interaction roles of the
Realtime-SFT corpus (\S\ref{sec:realtime-sft})---a standing
instruction that must fire exactly when its condition is met, and
continuous real-time commentary---under real-world timing.
